\documentclass[11pt, a4paper]{article}

% --- Geometry & Page Setup ---
\usepackage[a4paper, margin=2.3cm]{geometry}

% --- Font & Language Configuration ---
\usepackage[T1]{fontenc}
\usepackage[utf8]{inputenc}
\usepackage[english]{babel}
\usepackage{helvet}
\renewcommand{\familydefault}{\sfdefault}

% --- Essential Packages ---
\usepackage{amsmath, amssymb}
\usepackage{booktabs}
\usepackage{xcolor}
\usepackage{listings}
\usepackage{enumitem}
\usepackage{hyperref}

% --- Custom Color Definitions ---
\definecolor{brandblue}{RGB}{30, 58, 138}
\definecolor{brandaccent}{RGB}{13, 148, 136}
\definecolor{codebg}{RGB}{248, 250, 252}
\definecolor{codeframe}{RGB}{226, 232, 240}
\definecolor{commentgray}{RGB}{100, 116, 139}

% --- Setup Hyperlinks ---
\hypersetup{
    colorlinks=true,
    linkcolor=brandblue,
    citecolor=brandaccent,
    urlcolor=brandblue,
    pdftitle={Assignment 2 Detailed Report - Dynamics and Energies},
    pdfauthor={Filippo Segalla}
}

% --- Setup Listings (Code Environment) ---
\lstset{
    backgroundcolor=\color{codebg},
    basicstyle=\ttfamily\footnotesize,
    breaklines=true,
    captionpos=b,
    commentstyle=\color{commentgray},
    frame=single,
    rulecolor=\color{codeframe},
    keywordstyle=\color{brandblue}\bfseries,
    stringstyle=\color{brandaccent},
    showstringspaces=false,
    tabsize=2
}

% --- Custom Section Title Formatting ---
\makeatletter
\renewcommand{\section}{\@startsection{section}{1}{\z@}%
    {-3.5ex \@plus -1ex \@minus -.2ex}%
    {2.3ex \@plus .2ex}%
    {\normalfont\Large\bfseries\color{brandblue}}}
\renewcommand{\subsection}{\@startsection{subsection}{2}{\z@}%
    {-3.25ex \@plus -1ex \@minus -.2ex}%
    {1.5ex \@plus .2ex}%
    {\normalfont\large\bfseries\color{brandaccent}}}
\makeatother

% --- Document Metadata ---
\title{
    \vspace{-1cm}
    \color{brandblue}\textbf{Advanced Control Systems}\\[0.3em]
    \Large\textbf{Assignment 2: Kinetic and Potential Energy Computation, CoM Positions, Partial Jacobians, and Inertia Matrix of a PRR Manipulator}
}
\author{\textbf{Student Name:} Filippo Segalla \quad|\quad \textbf{Student ID:} VR538212\\[0.2em]
\small Department of Engineering for Innovative Medicine (DIMI) }
\date{\today}

\begin{document}

\maketitle
\thispagestyle{empty}

\begin{abstract}
\noindent This report presents the dynamic modeling and energy characterization of a 3-DOF PRR (Prismatic-Revolute-Revolute) robotic manipulator defined by the URDF specification. The study details the physical rationale and analytical derivation behind every dynamic modeling phase: link Center of Mass (CoM) position vectors, partial translational and rotational Jacobians ($\mathbf{J}_P^{l_i}, \mathbf{J}_O^{l_i}$), link local inertia tensors, the configuration-dependent generalized inertia matrix $\mathbf{B}(q)$, total kinetic energy $\mathcal{T}(q, \dot{q})$, and gravitational potential energy $\mathcal{U}(q)$. Every step is derived symbolically in parametric form and evaluated numerically using MATLAB.
\end{abstract}

\hrule
\vspace{1.5em}

\section{System Architecture \& Dynamic Parameters}

\subsection{Description of the PRR Manipulator Dynamic Model}
\paragraph{What is being done:}
We formulate the dynamic energy expressions for a spatial 3-degree-of-freedom PRR manipulator. The physical properties of the links are extracted directly from the URDF file (\texttt{00001\_PRR\_zzx} \texttt{FILIPPO} \texttt{SEGALLA.urdf}):
\begin{enumerate}[leftmargin=1.5em]
    \item \textbf{Link 1 (Prismatic Joint, $q_1$):} A rectangular box ($0.036\text{ m} \times 0.036\text{ m} \times 0.5\text{ m}$) with mass $m_1 = 1.0\text{ kg}$.
    \item \textbf{Link 2 (Revolute Joint, $q_2$):} A solid cylinder ($r = 0.019\text{ m}, h = 0.4\text{ m}$) aligned along $Z_2$ with mass $m_2 = 1.0\text{ kg}$.
    \item \textbf{Link 3 (Revolute Joint, $q_3$):} A solid cylinder ($r = 0.019\text{ m}, h = 0.4\text{ m}$) aligned along $X_3$ with mass $m_3 = 1.0\text{ kg}$.
\end{enumerate}

\paragraph{Why it is done:}
A complete dynamic description requires combining geometric kinematics with inertial properties (masses and inertia tensors). Establishing exact mass values and link geometry formulas is necessary to set up the Lagrangian $\mathcal{L} = \mathcal{T} - \mathcal{U}$, which governs the physical equations of motion.

\subsection{Local Inertia Tensors (\texorpdfstring{$\mathbf{I}_{l_i}^i$}{I\_li\_i})}
\paragraph{What is being done:}
For each link, the central inertia tensor $\mathbf{I}_{l_i}^i$ is evaluated with respect to a frame centered at its Center of Mass (CoM) and aligned with the link's principal axes of inertia:
\begin{itemize}[leftmargin=1.5em]
    \item \textbf{Link 1 (Parallelepiped Box):}
    \begin{equation}
        I_{1xx} = I_{1yy} = \frac{1}{12} m_1 (0.036^2 + 0.5^2), \quad I_{1zz} = \frac{1}{12} m_1 (0.036^2 + 0.036^2)
    \end{equation}
    \item \textbf{Link 2 (Cylinder along $Z_2$):}
    \begin{equation}
        I_{2xx} = I_{2yy} = \frac{1}{12} m_2 (3 \cdot 0.019^2 + 0.4^2), \quad I_{2zz} = \frac{1}{2} m_2 (0.019^2)
    \end{equation}
    \item \textbf{Link 3 (Cylinder along $X_3$):}
    \begin{equation}
        I_{3xx} = \frac{1}{2} m_3 (0.019^2), \quad I_{3yy} = I_{3zz} = \frac{1}{12} m_3 (3 \cdot 0.019^2 + 0.4^2)
    \end{equation}
\end{itemize}

\paragraph{Why it is done:}
Inertia tensors quantify a body's resistance to rotational acceleration. By expressing inertia tensors in local central frames aligned with the principal axes, off-diagonal cross-products vanish ($\mathbf{I}_{l_i}^i = \text{diag}(I_{ixx}, I_{iyy}, I_{izz})$). When the robot moves, these constant local tensors are transformed into the global frame using rotation matrices: $\mathbf{I}_{l_i}(q) = \mathbf{R}_i(q) \mathbf{I}_{l_i}^i \mathbf{R}_i^T(q)$.

\section{Centers of Mass \& Partial Jacobians}

\subsection{Centers of Mass (CoM) Positions in Base Frame}
\paragraph{What is being done:}
We locate the position vector $\mathbf{p}_{l_i}$ of each link's CoM relative to the global base frame $\Sigma_0$ using cumulative homogeneous transformation matrices:
\begin{equation}
    \begin{bmatrix} \mathbf{p}_{l_i} \\ 1 \end{bmatrix} = \mathbf{T}_i^0(q) \begin{bmatrix} \mathbf{p}_{l_i}^i \\ 1 \end{bmatrix}
\end{equation}

Where the local CoM vectors relative to the D-H origins are:
\begin{equation}
    \mathbf{p}_{l_1}^1 = \begin{bmatrix} 0 \\ 0 \\ -0.25 \\ 1 \end{bmatrix}, \quad
    \mathbf{p}_{l_2}^2 = \begin{bmatrix} 0 \\ 0 \\ -0.20 \\ 1 \end{bmatrix}, \quad
    \mathbf{p}_{l_3}^3 = \begin{bmatrix} -0.20 \\ 0 \\ 0 \\ 1 \end{bmatrix}
\end{equation}

Evaluating the transformations symbolically yields:
\begin{equation}
    \mathbf{p}_{l_1}(q) = \begin{bmatrix} 0 \\ 0 \\ q_1 + 0.40 \end{bmatrix}, \quad
    \mathbf{p}_{l_2}(q) = \begin{bmatrix} 0 \\ 0 \\ q_1 + 0.85 \end{bmatrix}, \quad
    \mathbf{p}_{l_3}(q) = \begin{bmatrix} -0.20 \sin q_2 \cos q_3 \\ 0.20 \cos q_2 \cos q_3 \\ q_1 + 0.20 \sin q_3 + 1.05 \end{bmatrix}
\end{equation}

\paragraph{Why it is done:}
König's theorem states that the total kinetic energy of a rigid link is the sum of translational energy at its CoM and rotational energy about its CoM. Knowing $\mathbf{p}_{l_i}(q)$ is necessary for calculating both the linear velocities of the link centers of mass and the gravitational potential energy $\mathcal{U}_i = -m_i \mathbf{g}_0^T \mathbf{p}_{l_i}$.

\subsection{Partial Jacobians (\texorpdfstring{$\mathbf{J}_P^{l_i}, \mathbf{J}_O^{l_i}$}{JP\_li, JO\_li})}
\paragraph{What is being done:}
Partial Jacobians map generalized joint velocities $\dot{q}$ to the linear velocity $\dot{\mathbf{p}}_{l_i}$ and angular velocity $\boldsymbol{\omega}_i$ of link $i$:
\begin{equation}
    \dot{\mathbf{p}}_{l_i} = \mathbf{J}_P^{l_i}(q) \dot{q}, \quad \boldsymbol{\omega}_i = \mathbf{J}_O^{l_i}(q) \dot{q}
\end{equation}

The columns $j$ ($j \le i$) of the partial Jacobians are computed as:
\begin{itemize}[leftmargin=1.5em]
    \item \textbf{Prismatic Joint $j$:} $\mathbf{J}_{Pj}^{l_i} = \mathbf{z}_{j-1}$, $\mathbf{J}_{Oj}^{l_i} = \mathbf{0}$.
    \item \textbf{Revolute Joint $j$:} $\mathbf{J}_{Pj}^{l_i} = \mathbf{z}_{j-1} \times (\mathbf{p}_{l_i} - \mathbf{p}_{j-1})$, $\mathbf{J}_{Oj}^{l_i} = \mathbf{z}_{j-1}$.
    \item \textbf{Inactive Joints ($j > i$):} $\mathbf{J}_{Pj}^{l_i} = \mathbf{0}$, $\mathbf{J}_{Oj}^{l_i} = \mathbf{0}$.
\end{itemize}

\paragraph{Why it is done:}
Since link $i$ is affected only by joints $1, \dots, i$, columns beyond $i$ are zero. Partial Jacobians allow us to write the kinetic energy of each link in a compact quadratic form without explicitly dealing with velocity vector products.

\section{Generalized Inertia Matrix \texorpdfstring{$\mathbf{B}(q)$}{B(q)}}

\subsection{Mathematical Formulation}
\paragraph{What is being done:}
The generalized inertia matrix $\mathbf{B}(q) \in \mathbb{R}^{3 \times 3}$ is constructed by summing the translational and rotational kinetic energy contributions of all links:
\begin{equation}
    \mathbf{B}(q) = \sum_{i=1}^{3} \left( m_i (\mathbf{J}_P^{l_i})^T \mathbf{J}_P^{l_i} + (\mathbf{J}_O^{l_i})^T \mathbf{R}_i \mathbf{I}_{l_i}^i \mathbf{R}_i^T \mathbf{J}_O^{l_i} \right)
\end{equation}

Evaluating this sum symbolically yields the parametric expression:
\begin{equation}
\setlength{\arraycolsep}{3pt}
\mathbf{B}(q) = \begin{bmatrix}
m_1 + m_2 + m_3 & 0 & \frac{m_3}{5} \cos q_3 \\[0.5em]
0 & I_{2yy} + I_{3yy}\cos^2 q_3 + I_{3xx}\sin^2 q_3 + \frac{m_3}{25}\cos^2 q_3 & 0 \\[0.5em]
\frac{m_3}{5} \cos q_3 & 0 & I_{3zz} + \frac{m_3}{25}
\end{bmatrix}
\end{equation}

Substituting numerical parameters gives:
\begin{equation}
\mathbf{B}(q) = \begin{bmatrix}
3.0 & 0 & 0.2 \cos q_3 \\
0 & 0.1068 - 0.05324 \sin^2 q_3 & 0 \\
0.2 \cos q_3 & 0 & 0.05342
\end{bmatrix}
\end{equation}

\paragraph{Why it is done \& Physical Properties:}
\begin{enumerate}[leftmargin=1.5em]
    \item \textbf{Symmetry ($\mathbf{B} = \mathbf{B}^T$):} Reflects the reciprocal dynamic coupling between joints.
    \item \textbf{Positive Definiteness ($\mathbf{B} > 0$):} Ensures kinetic energy is strictly positive for non-zero velocities ($\mathcal{T} > 0$ for $\dot{q} \neq 0$).
    \item \textbf{Coupling Term ($B_{13} = 0.2 \cos q_3$):} Quantifies the inertial interaction between vertical base translation ($q_1$) and wrist rotation ($q_3$).
    \item \textbf{Invariance to $q_2$:} $\mathbf{B}(q)$ is independent of $q_2$ because rotation about the vertical $Z$-axis does not alter the manipulator's mass distribution relative to gravity or rotation axes.
\end{enumerate}

\section{Kinetic \& Potential Energy Computation}

\subsection{Total Kinetic Energy \texorpdfstring{$\mathcal{T}(q, \dot{q})$}{T(q, dq)}}
\paragraph{What is being done:}
The total kinetic energy is formulated as a quadratic expression in joint velocities:
\begin{equation}
    \mathcal{T}(q, \dot{q}) = \frac{1}{2} \dot{q}^T \mathbf{B}(q) \dot{q}
\end{equation}

Evaluating the matrix product numerically yields:
\begin{equation}
    \mathcal{T}(q, \dot{q}) = 1.5 \dot{q}_1^2 + 0.2 \dot{q}_1 \dot{q}_3 \cos q_3 - 0.02662 \dot{q}_2^2 \sin^2 q_3 + 0.05342 \dot{q}_2^2 + 0.02671 \dot{q}_3^2
\end{equation}

\paragraph{Why it is done:}
Kinetic energy stores the mechanical work done by joint actuators during motion. The quadratic term $\frac{1}{2} B_{ii} \dot{q}_i^2$ represents the effective inertia at joint $i$, while cross-terms $B_{ij} \dot{q}_i \dot{q}_j$ represent dynamic cross-coupling.

\subsection{Total Potential Energy \texorpdfstring{$\mathcal{U}(q)$}{U(q)}}
\paragraph{What is being done:}
Potential energy due to gravity ($\mathbf{g}_0 = [0, 0, -g]^T$) is calculated as:
\begin{equation}
    \mathcal{U}(q) = -\sum_{i=1}^{3} m_i \mathbf{g}_0^T \mathbf{p}_{l_i}(q)
\end{equation}

Evaluating symbolically gives:
\begin{equation}
    \mathcal{U}(q) = g m_1 \left(q_1 + \frac{2}{5}\right) + g m_2 \left(q_1 + \frac{17}{20}\right) + g m_3 \left(q_1 + \frac{\sin q_3}{5} + \frac{21}{20}\right)
\end{equation}

Substituting numerical values ($g = 9.81\text{ m/s}^2, m_i = 1\text{ kg}$):
\begin{equation}
    \mathcal{U}(q) = 29.43 q_1 + 1.962 \sin q_3 + 22.56
\end{equation}

\paragraph{Why it is done:}
Potential energy depends strictly on position coordinates. For this PRR robot, $\mathcal{U}(q)$ depends on $q_1$ (vertical translation) and $q_3$ (elevation of Link 3), but is completely independent of $q_2$, since rotation about $Z_0$ moves the robot within a horizontal plane perpendicular to gravity.

\section{MATLAB Implementation \& Console Outputs}

\subsection{Overview of the Script (\texttt{SecondAssignment.m})}
\paragraph{What is being done:}
The script \texttt{SecondAssignment.m} implements the dynamic energy pipeline. It first computes all dynamic matrices symbolically using generic letters ($m_i, I_{ixx}, g$), displays the parametric expressions, and then performs numerical substitution to display clean decimal results using \texttt{vpa(..., 4)}.

\paragraph{Why it is done:}
Displaying both symbolic parametric formulas and evaluated numerical matrices allows verification of the underlying physical equations while providing concrete values for simulation and controller design.

\subsection{Verified Console Logs}
Below are the exact console outputs generated by executing \texttt{SecondAssignment.m} in MATLAB:

\begin{lstlisting}[caption={MATLAB Console Output: Part 1 - Purely Symbolic Formulas (Letters).}]
------------------------------------------------------------
PURELY SYMBOLIC FORMULAS (LETTERS)
------------------------------------------------------------

--- CoM Positions p_l in Base Frame Sigma_0 ---
p_l1 =
          0
          0
 q1 + 2/5

p_l2 =
           0
           0
 q1 + 17/20

p_l3 =
 -(m3*cos(q3)*sin(q2))/5
  (m3*cos(q3)*cos(q2))/5
 q1 + sin(q3)/5 + 21/20

--- Symbolic Inertia Matrix B(q) (3x3) ---
[ m1 + m2 + m3,                                                 0,                            (m3*cos(q3))/5]
[            0, I2yy + I3yy*cos(q3)^2 + I3xx*sin(q3)^2 + (m3*cos(q3)^2)/25,                                         0]
[ (m3*cos(q3))/5,                                                 0, I3zz + m3/25]

--- Symbolic Kinetic Energy T(q, dq) ---
(dq1*(dq1*(m1 + m2 + m3) + (dq3*m3*cos(q3))/5))/2 + (dq3*((dq1*m3*cos(q3))/5 + dq3*(I3zz + m3/25)))/2 + (dq2^2*(I2yy + I3yy*cos(q3)^2 + I3xx*sin(q3)^2 + (m3*cos(q3)^2)/25))/2

--- Symbolic Potential Energy U(q) ---
g*m1*(q1 + 2/5) + g*m2*(q1 + 17/20) + g*m3*(q1 + sin(q3)/5 + 21/20)
\end{lstlisting}

\begin{lstlisting}[caption={MATLAB Console Output: Part 2 - Evaluated Numerical Values (Decimals).}]
------------------------------------------------------------
EVALUATED NUMERICAL VALUES (DECIMALS)
------------------------------------------------------------

--- Evaluated Inertia Matrix B(q) (3x3) ---
[        3.0,                               0,    0.2*cos(q3)]
[          0,  0.1068 - 0.05324*sin(q3)^2,              0]
[ 0.2*cos(q3),                               0,       0.05342]

--- Evaluated Kinetic Energy T(q, dq) ---
1.5*dq1^2 + 0.2*dq1*dq3*cos(q3) - 0.02662*dq2^2*sin(q3)^2 + 0.05342*dq2^2 + 0.02671*dq3^2

--- Evaluated Potential Energy U(q) (with g = 9.81) ---
29.43*q1 + 1.962*sin(q3) + 22.56
============================================================
\end{lstlisting}

\end{document}